\section{Ondas eletromagnéticas}

A lei de Faraday nos diz que a variação no campo magnéticao gera um campo elétrico, enquanto a lei de Ampère-Maxwell nos diz que variação no campo elétrico gera um campo magnético.
Portanto, um uma perturbação eletromagnética gera um pulso que se propaga pelo espaço.
Nessa sessão, será mostrado que essa propagação dos campos $\mathbf E$ e $\mathbf B$ obedece uma equação de onda da forma
\begin{equation}
    v^2 \nabla^2 \mathbf F = \partial_t^2 \mathbf F \;,
    \label{eq: wave}
\end{equation}
\noindent onde $v$ é a velocidade de propação da onda.
Vamos considerar as equações de Maxwell no vácuo e na ausência de cargas e correntes elétricas ($\rho = 0,\, \mathbf j = 0$).
Pela identidade \ref{eq: nabla_identity}, temos
\begin{equation}
    \nabla \times \nabla \times \mathbf E = \nabla (\nabla \cdot \mathbf E) - \nabla^2 \mathbf E = -\nabla^2 \mathbf E \;,
    \label{eq: wave_E_pt1}
\end{equation}
\noindent onde foi utilizada a eq. de Poisson \eqref{eq: poisson} com $\rho = 0$. Por outro lado, temos
\begin{align}
    \nabla \times \nabla \times \mathbf E &= \nabla \times (-\partial_t \mathbf B) = -\partial_t (\nabla \times \mathbf B) \notag \\
    &= -\partial_t (\epsilon_0\mu_0\partial_t \mathbf E) = -\epsilon_0\mu_0\partial_t^2 \mathbf E
    \label{eq: wave_E_pt2}
\end{align}
\noindent onde foram utilizadas a lei de Faraday \eqref{eq: faraday} e a lei de Ampère-Maxwell \eqref{eq: dif_amperemaxwell} com $\mathbf j = 0$.
Juntando as eqs.\eqref{eq: wave_E_pt1} e \eqref{eq: wave_E_pt2}, obtemos
\begin{equation}
    \nabla^2 \mathbf E = \epsilon_0\mu_0\partial_t^2 \mathbf E \;.
    \label{eq: wave_E}
\end{equation}

Analogamente, aplicando a identidade \ref{eq: nabla_identity} ao campo $\mathbf B$ e utilizando a lei de Gauss para o magnetismo \eqref{eq: dif_gauss_mag}, a lei Ampère-Maxwell \eqref{eq: dif_amperemaxwell} com $\mathbf j = 0$, e a lei de Faraday \eqref{eq: faraday}, obtemos
\begin{equation}
    \nabla^2 \mathbf B = \epsilon_0\mu_0\partial_t^2 \mathbf B \;.
    \label{eq: wave_B}
\end{equation}

Assim, os dois campos possuem a forma da eq.\eqref{eq: wave}.
Concluímos que as perturbações se propagam na forma de ondas eletromagnéticas, cuja velocidade é $c \equiv (\epsilon_0 \mu_0)^{-1/2} = 299792458$ m/s.

É possível mostrar que soluções das equações de Maxwell são da forma
\begin{equation}
    \mathbf B = \nabla \times \mathbf A \;, \quad \mathbf E = -\nabla \phi - \partial_t \mathbf A \;,
    \label{eq: maxwell_solutions}
\end{equation}
\noindent onde $\phi$ e $\mathbf A$ são respectivamente o potencial escalar e o potencial vetor, que obedecem a relação
\begin{equation}
    \nabla \cdot \mathbf A + \epsilon_0\mu_0\partial_t \phi /c = 0 \;.
    \label{eq: lorentz_condition}
\end{equation}
`
